\scanpage{59}
\mspage{58}

$\varphi(L)$ est la plus grande sous-figure commune de $\varphi(F),\varphi(G)$. L'hyp. $\varphi(F)\unsure{\gtrless|\lessgtr|\gtreqless|\lesseqgtr}\varphi(G)$ équivaut à
\[
\begin{aligned}
&|\varphi(F)|\cap|\varphi(G)|=|\varphi(L)|\qquad\text{i.e.}\\
&\operatorname{omb}(F)\cap\operatorname{omb}(G)=\operatorname{omb}(L),
\end{aligned}
\]
mais c'est (en vertu de \unsure{AtCL (ap a)|AtCL (op a)}) la condition pour que $F\gtrless G$. OK.

Prouvons b), i.e. soient $F,G$ avec $\varphi(F)\ll\varphi(G)$, prouvons $F\ll G$. \deleted{\unsure{Il}} Utilisant a) déjà prouvé, on est ramené au cas $X,Y\in M$.

Mais c'est encore faux!

\uline{Ex} \deleted{\deleted{$X$ espace affine sur $\unsure{\mathbb{R}|R}$. On prend $M\unsure{=|-}$ simplexes affines, $\leq$ \unsure{pour} sous-simplexe,} \inserted{On reprend $\unsure{M}$ comme dans l'exemple précédent,} mais en renforçant cette fois la relation $\ll$, pour la rendre conforme à la version ordinaire (``\unsure{simpliciale}'' ou ``polyédrale''\unsure{.})}

Soit $L$ un ens., considérons l'ens $\unsure{\mathbb{F}|\mathrm{F}}\mathrm{ig}\,\mathrm{pol}(L)$ des figures \uline{polyédrales} dans $L$. On y prend $\leq$ et $\ll$ comme d'habitude\inserted{, dans le cas polyédral}, étant entendu que $\ll$ signifie $\overset{\mathrm{pol}}{\ll}$, tandis qu'on prend la compatibilité $\gtrless$ \uline{\unsure{ens}} (\unsure{que|par} \unsure{sur|sous} ensembles \unsure{en|ou} polyèdres). (D'avoir renforcé le $\underset{\mathrm{ens}}{\ll}$ en $\underset{\mathrm{pol}}{\ll}$ ne modifie pas la validité des conditions $\unsure{\mathrm{Prat}|\mathrm{Prot}}\,1'-\mathrm{Prat}\,6'$.) Considérons deux \deleted{figures}\inserted{str. de} polyèdres (p. ex. de \deleted{\unsure{qui sont}} type $\Delta_1,\unsure{\Delta_n|\Delta_1}$) $X,Y$, avec $X\underset{\mathrm{ens}}{\ll}Y$ et non $X\underset{\mathrm{pol}}{\ll}Y$, on a \unsure{le|là} contre-exemple.
